\section{هفته‌ی 9: مدل‌سازی بیزی 1: پیشین، پسین و تقریب شبکه‌ای}

\begin{quote}
\textbf{فصل 4: مدل‌سازی بیزی} · درس 1 ساعت · \lr{lab} 2 ساعت
\textbf{داده‌ی اصلی:} یک نسبت یادآوری (30 از 50) و شرکت‌کننده‌ی \lr{SDT}
هفته‌ی 2 (68/100 \lr{hit}، 11/100 \lr{false alarm).}
\end{quote}

\begin{center}\rule{0.5\linewidth}{0.5pt}\end{center}

\subsection{واژه‌های کلیدی}

واژه‌های مهم این درس با معادل فارسی‌شان (حتی اگر در درس‌های قبل آمده
باشند).

\begin{longtable}[]{@{}
  >{\raggedright\arraybackslash}p{(\linewidth - 4\tabcolsep) * \real{0.3333}}
  >{\raggedright\arraybackslash}p{(\linewidth - 4\tabcolsep) * \real{0.3333}}
  >{\raggedright\arraybackslash}p{(\linewidth - 4\tabcolsep) * \real{0.3333}}@{}}
\caption{\centering واژه‌های کلیدی}\tabularnewline
\toprule\noalign{}
\begin{minipage}[b]{\linewidth}\raggedright
\lr{English}
\end{minipage} & \begin{minipage}[b]{\linewidth}\raggedright
فارسی
\end{minipage} & \begin{minipage}[b]{\linewidth}\raggedright
توضیح کوتاه
\end{minipage} \\
\midrule\noalign{}
\endfirsthead
\toprule\noalign{}
\begin{minipage}[b]{\linewidth}\raggedright
\lr{English}
\end{minipage} & \begin{minipage}[b]{\linewidth}\raggedright
فارسی
\end{minipage} & \begin{minipage}[b]{\linewidth}\raggedright
توضیح کوتاه
\end{minipage} \\
\midrule\noalign{}
\endhead
\bottomrule\noalign{}
\endlastfoot
\lr{prior} & توزیع پیشین \lr{(prior)} & باور درباره‌ی پارامتر پیش از دیدن
داده \\
\lr{posterior} & توزیع پسین \lr{(posterior)} & باور درباره‌ی پارامتر بعد
از دیدن داده \\
\lr{grid approximation} & تقریب شبکه‌ای \lr{(grid)} & محاسبه‌ی پسین روی
شبکه‌ای از مقادیر \\
\lr{signal detection theory (SDT)} & نظریه‌ی تشخیص علامت & مدلی که حساسیت
را از سوگیری پاسخ جدا می‌کند \\
\lr{hit} & اصابت \lr{(hit)} & پاسخ «بله» وقتی علامت وجود دارد \\
\lr{Bayes' rule} & قاعده‌ی بیز & پسین \ensuremath{\propto} درست‌نمایی × پیشین \\
\lr{parameter} & پارامتر & کمیت قابل تنظیم مدل که بین افراد یا شرایط فرق
می‌کند \\
\lr{likelihood} & درست‌نمایی & احتمال داده به‌عنوان تابعی از پارامترها \\
\lr{evidence} & شواهد & اطلاعاتی که برای تصمیم جمع می‌شود \\
\lr{beta-binomial model} & مدل \lr{beta-binomial} & پیشین \lr{Beta} با
داده‌ی دوجمله‌ای \\
\lr{binomial distribution} & توزیع دوجمله‌ای & توزیع تعداد موفقیت در
\lr{n} کوشش \\
\lr{MAP estimate} & برآورد \lr{MAP} & مُد توزیع پسین \\
\lr{credible interval} / \lr{HDI} & بازه‌ی باورپذیر \lr{(credible
interval)} & بازه‌ای که با احتمال مشخص پارامتر را دارد \\
\lr{confidence interval (CI)} & بازه‌ی اطمینان & بازه‌ای که با روش مشخص
پارامتر واقعی را در بر می‌گیرد \\
\lr{prior predictive check} & بررسی پیش‌بینی پیشین & شبیه‌سازی داده از
پیشین برای سنجش معقول بودن آن \\
\lr{simulation} & شبیه‌سازی & تولید داده‌ی مصنوعی با اجرای مدل \\
\lr{sensitivity (d′)} & حساسیت & فاصله‌ی توزیع علامت و نویز \\
\lr{estimate} / \lr{estimator} & برآورد / برآوردگر & مقدار به‌دست‌آمده
برای پارامتر / روش به دست آوردن آن \\
\lr{maximum likelihood estimate (MLE)} & برآورد بیشینه‌ی درست‌نمایی &
پارامترهایی که داده را محتمل‌ترین می‌کنند \\
\lr{cognitive model} & مدل شناختی & توصیف دقیق و ریاضی یا محاسباتی از یک
فرایند ذهنی \\
\lr{trial} & کوشش \lr{(trial)} & یک بار انجام تکلیف توسط شرکت‌کننده \\
\lr{conjugate prior} & پیشین مزدوج \lr{(conjugate)} & پیشینی که پسین
هم‌خانواده می‌دهد \\
\lr{forgetting curve} & منحنی فراموشی & کاهش یادآوری با گذشت زمان \\
\lr{accuracy} & دقت & نسبت پاسخ‌های درست \\
\end{longtable}

\begin{center}\rule{0.5\linewidth}{0.5pt}\end{center}

\subsection{اهداف یادگیری}

در پایان این هفته، دانشجویان می‌توانند:

\begin{enumerate}
\def\labelenumi{\arabic{enumi}.}
\tightlist
\item
  قاعده‌ی بیز را برای پارامترها \textbf{بیان کنند} و نقش پیشین
  \lr{(prior)}، درست‌نمایی، شواهد \lr{(evidence)} و پسین \lr{(posterior)}
  را \textbf{مشخص کنند}.
\item
  پسین را برای مدل \lr{beta-binomial} به‌صورت تحلیلی و با تقریب شبکه‌ای
  به‌صورت عددی، برای یک و دو پارامتر، \textbf{حساب کنند}.
\item
  یک پسین را \textbf{خلاصه کنند} (میانگین، میانه، \lr{MAP}، انحراف
  معیار، بازه‌های دم‌برابر و بیشترین چگالی) و \lr{credible interval} را از
  بازه‌ی اطمینان \textbf{متمایز کنند}.
\item
  \textbf{توضیح دهند} که پیشین و مقدار داده با هم چطور پسین را تعیین
  می‌کنند، و با شبیه‌سازی پیش‌بینی پیشین \textbf{بررسی کنند} که پیشین معقول
  است یا نه.
\item
  احتمال‌های پسین فرضیه‌های معنادار روان‌شناختی را \textbf{تفسیر کنند}
  (مثلاً «\(P(d' > 1 \mid \text{\lr{data}})\)»).
\end{enumerate}

\begin{center}\rule{0.5\linewidth}{0.5pt}\end{center}

\subsection{مقدمه}

\begin{quote}
\emph{«با توجه به این داده‌ها، با احتمال 95\% حساسیت این شرکت‌کننده بین
1.2 و 2.1 است.»} \emph{«احتمال اینکه دست‌کاری پاداش شرکت‌کننده‌ها را
\lr{liberal}تر کرده باشد 0.98 است.»}
\end{quote}

بازه‌های اطمینان فراوانی‌گرا اجازه‌ی چنین جمله‌هایی را \textbf{نمی‌دهند}
(هفته‌ی 5). استنباط بیزی می‌دهد، به یک قیمت: باید بگوییم \textbf{پیش از}
دیدن داده درباره‌ی پارامترها چه باوری داشتیم، به شکل یک \textbf{توزیع
پیشین}. بعد قاعده‌ی بیز پیشین را با درست‌نمایی (که بلدیم بنویسیم) ترکیب
می‌کند و یک \textbf{توزیع پسین} می‌سازد: همه‌ی آنچه مدل و داده با هم
درباره‌ی پارامترها می‌گویند.

روش‌های بیزی چند دردسر قبلی ما را هم آسان‌تر می‌کنند: عدم قطعیت مجانی به
دست می‌آید (بدون \lr{Hessian} یا \lr{bootstrap)}، مدل‌های سلسله‌مراتبی چند
شرکت‌کننده طبیعی می‌شوند (هفته‌ی 13)، و مقایسه‌ی مدل‌ها یک جواب اصولی پیدا
می‌کند (هفته‌ی 10). هزینه‌اش محاسبه است، و برای همین هفته‌های 11 تا 13
درباره‌ی نمونه‌گیری‌اند.

\begin{quote}
\textbf{یادآوری (هفته‌ی 1):} این هفته درباره‌ی \textbf{تحلیل داده‌ی بیزی}
است؛ روشی که ما پژوهشگران برای برآورد پارامترهای \emph{هر} مدلی به کار
می‌بریم. هیچ ادعایی ندارد که ذهن بیزی است. \emph{مدل‌های بیزی شناخت} در
هفته‌ی 15 می‌آیند.
\end{quote}

\paragraph{این موضوع کجا به کار می‌رود؟}

\begin{itemize}
\tightlist
\item
  \textbf{نمونه‌های کوچک و موردهای تکی.} در عصب‌روان‌شناسی بالینی اغلب یک
  بیمار و چند ده \lr{trial} داریم. تحلیل بیزی با یک \lr{prior} معقول،
  عدم قطعیت صادقانه و قابل استفاده‌ای می‌دهد، جایی که \lr{MLE} ناپایدار
  است.
\item
  \textbf{گزاره‌های احتمالی مستقیم} که پزشکان و سیاست‌گذاران می‌فهمند: «با
  احتمال 97\% آموزش جدید حافظه را بهتر کرده».
\item
  \textbf{آزمون‌های تطبیقی \lr{(adaptive testing):}} آزمون‌های کامپیوتری
  (مثلاً آزمون‌های بینایی و شنوایی، امتحان‌های تطبیقی) بعد از هر پاسخ
  \lr{posterior} را به‌روز می‌کنند و آیتم بعدی را جایی انتخاب می‌کنند که
  بیشترین اطلاعات را می‌دهد.
\item
  \textbf{ترکیب شواهد:} \lr{posterior} امروز \lr{prior} فرداست، پس نتایج
  مطالعه‌های قبلی را می‌شود به‌طور رسمی وارد تحلیل کرد.
\item
  بیشتر مدل‌های سلسله‌مراتبی در مقاله‌های اخیر به همین روش برازش می‌شوند
  (هفته‌ی 13).
\end{itemize}

\begin{center}\rule{0.5\linewidth}{0.5pt}\end{center}

\Needspace{11\baselineskip}

\subsection{جزوه‌ی درس}

\paragraph{به زبان ساده}

تا اینجا \emph{یک} بهترین مقدار پارامتر پیدا می‌کردیم و بعد می‌پرسیدیم
چقدر ممکن است تغییر کند. روش بیزی در عوض دانش ما درباره‌ی پارامتر را به
شکل یک \textbf{توزیع احتمال} کامل توصیف می‌کند. بعد داده‌ها به هر مقدار
ممکن به اندازه‌ی درست‌نمایی‌اش (هفته‌ی 4) وزن تازه می‌دهند. نتیجه،
\textbf{\lr{posterior} (پسین)}، می‌گوید هر مقدار \emph{بعد} از دیدن داده
چقدر معقول است. به زبان ساده: \emph{\lr{posterior} \ensuremath{\propto} \lr{likelihood} ×
\lr{prior}}. از \lr{posterior} می‌توانیم یک بهترین حدس (میانگین آن)، یک
عدم قطعیت (پراکندگی آن) و گزاره‌های مستقیمی مثل «احتمال اینکه \(d' > 1\)
باشد 0.93 است» بخوانیم. با داده‌ی کم \lr{prior} اثر زیادی دارد؛ با داده‌ی
زیاد \lr{posterior} زیر سلطه‌ی درست‌نمایی است و به \lr{MLE} نزدیک می‌شود.
محاسبه‌ی \lr{posterior} برای یک یا دو پارامتر آسان است (روی \lr{grid)} و
برای بیشتر از آن به روش‌های نمونه‌گیری نیاز دارد (هفته‌های 11 تا 13).

\textbf{زمان‌بندی پیشنهادی (60 دقیقه):} قاعده‌ی بیز 10 · پیشین‌ها 10 ·
\lr{beta-binomial} 12 · تقریب شبکه‌ای 10 · خلاصه‌ها و \lr{credible
interval}ها 12 · بررسی پیش‌بینی پیشین 6.

\subsubsection{1. قاعده‌ی بیز برای پارامترها}

\[
\underbrace{p(\boldsymbol\theta \mid \mathbf{y})}_{\text{\lr{posterior}}}
= \frac{\overbrace{p(\mathbf{y} \mid \boldsymbol\theta)}^{\text{\lr{likelihood}}}\;\overbrace{p(\boldsymbol\theta)}^{\text{\lr{prior}}}}{\underbrace{p(\mathbf{y})}_{\text{\lr{evidence (marginal likelihood)}}}},
\qquad
p(\mathbf{y}) = \int p(\mathbf{y}\mid\boldsymbol\theta)\,p(\boldsymbol\theta)\,d\boldsymbol\theta.
\]

(این هفته «\(\mathcal{M}\)» را در نمادگذاری نمی‌نویسیم: همه‌چیز مشروط به
مدل است.)

\begin{itemize}
\tightlist
\item
  \textbf{درست‌نمایی} دقیقاً همان تابع هفته‌ی 4 است.
\item
  \textbf{پیشین} آنچه را پیش از داده معلوم یا محتمل است بیان می‌کند.
\item
  \textbf{شواهد} \(p(\mathbf{y})\) به \(\boldsymbol\theta\) بستگی ندارد.
  برای برآورد، فقط ثابت نرمال‌ساز است که انتگرال پسین را 1 می‌کند:
\end{itemize}

\[
p(\boldsymbol\theta \mid \mathbf{y}) \propto p(\mathbf{y}\mid\boldsymbol\theta)\, p(\boldsymbol\theta).
\]

\textbf{پسین \ensuremath{\propto} درست‌نمایی × پیشین.} در هفته‌ی 10 شواهد ستاره‌ی مقایسه‌ی بیزی
مدل‌ها می‌شود.

\textbf{شهود.} پسین یک مصالحه است: جایی بالاست که مقادیر پارامتر هم
\emph{از پیش} محتمل‌اند و هم داده را محتمل می‌کنند. با داده‌ی کم پیشین مهم
است؛ با داده‌ی زیاد درست‌نمایی غالب می‌شود و پسین حول \lr{MLE} متمرکز
می‌شود.

\paragraph{خواندن ریاضی}

\begin{itemize}
\tightlist
\item
  \(p(\boldsymbol\theta \mid \mathbf{y})\): «احتمال (چگالی) مقدار
  پارامتر \(\boldsymbol\theta\) \emph{به شرط} داده». این همان چیزی است
  که می‌خواهیم.
\item
  \(p(\mathbf{y} \mid \boldsymbol\theta)\): درست‌نمایی، که به‌عنوان تابعی
  از \(\boldsymbol\theta\) خوانده می‌شود (هفته‌ی 4).
\item
  \(p(\boldsymbol\theta)\): \lr{prior} که پیش از داده انتخاب می‌شود.
\item
  \(p(\mathbf{y}) = \int p(\mathbf{y}\mid\boldsymbol\theta)p(\boldsymbol\theta)\,d\boldsymbol\theta\):
  انتگرال، «درست‌نمایی × \lr{prior»} را روی \emph{همه‌ی} مقدارهای ممکن
  پارامتر جمع می‌کند. یک عدد است و برای همه‌ی \(\boldsymbol\theta\)ها
  یکسان است، پس در برآورد فقط منحنی را طوری مقیاس می‌کند که مساحت کل 1
  شود.
\item
  \(\propto\) یعنی «متناسب با»: برابر، به‌جز یک ضریب ثابت. در عمل
  «درست‌نمایی × \lr{prior»} را در مقدارهای زیادی حساب می‌کنیم و بعد بر
  جمعشان تقسیم می‌کنیم، دقیقاً مثل جدول چهارنقطه‌ای مثال حل‌شده (0.0004479 /
  0.0009574 = 0.468).
\end{itemize}

\subsubsection{2. انتخاب پیشین‌ها}

\begin{longtable}[]{@{}
  >{\raggedright\arraybackslash}p{(\linewidth - 4\tabcolsep) * \real{0.3333}}
  >{\raggedright\arraybackslash}p{(\linewidth - 4\tabcolsep) * \real{0.3333}}
  >{\raggedright\arraybackslash}p{(\linewidth - 4\tabcolsep) * \real{0.3333}}@{}}
\caption{\centering 2. انتخاب پیشین‌ها}\tabularnewline
\toprule\noalign{}
\begin{minipage}[b]{\linewidth}\raggedright
نوع
\end{minipage} & \begin{minipage}[b]{\linewidth}\raggedright
مثال برای احتمال یادآوری \(\theta\)
\end{minipage} & \begin{minipage}[b]{\linewidth}\raggedright
چه وقت
\end{minipage} \\
\midrule\noalign{}
\endfirsthead
\toprule\noalign{}
\begin{minipage}[b]{\linewidth}\raggedright
نوع
\end{minipage} & \begin{minipage}[b]{\linewidth}\raggedright
مثال برای احتمال یادآوری \(\theta\)
\end{minipage} & \begin{minipage}[b]{\linewidth}\raggedright
چه وقت
\end{minipage} \\
\midrule\noalign{}
\endhead
\bottomrule\noalign{}
\endlastfoot
تخت / یکنواخت & \(\mathrm{Beta}(1, 1)\) & وقتی هر مقدار در دامنه واقعاً
به یک اندازه محتمل است (به‌ندرت درست است) \\
ضعیف‌آگاهی‌بخش \lr{(weakly informative)} & \(\mathrm{Beta}(2, 2)\): ترجیح
ملایم برای مقادیر میانی & \textbf{پیش‌فرض برای بیشتر مدل‌های شناختی}:
مقادیر پوچ را کنار می‌گذارد و بقیه را به داده می‌سپارد \\
آگاهی‌بخش \lr{(informative)} & \(\mathrm{Beta}(20, 5)\): «یادآوری بعد از
یک روز حدود 0.8 است» & وقتی مطالعات قبلی یا نظریه دانش واقعی می‌دهند \\
\end{longtable}

اصول: - پیشین‌ها باید \textbf{پیش از} دیدن داده تعیین شوند. - باید دامنه‌ی
پارامتر را رعایت کنند (مثلاً هیچ جرم پیشینی روی نرخ زوال منفی نباشد). -
یک پیشین «مبهم» روی یک مقیاس می‌تواند روی مقیاس دیگری خیلی آگاهی‌بخش باشد
(بخش \lr{D lab).} - پیشین‌هایتان را گزارش کنید و حساسیت نتیجه‌گیری‌ها را به
جایگزین‌های معقول بررسی کنید (هفته‌ی 10).

\subsubsection{\texorpdfstring{3. یک حالت با جواب بسته: مدل
\lr{beta-binomial}}{3. یک حالت با جواب بسته: مدل beta-binomial}}

برای \(y\) موفقیت در \(n\) \lr{trial} دوجمله‌ای با پیشین
\(\mathrm{Beta}(a_0, b_0)\) روی \(\theta\):

\[
p(\theta \mid y) \propto \underbrace{\theta^{y}(1-\theta)^{n-y}}_{\text{\lr{likelihood}}} \cdot \underbrace{\theta^{a_0 - 1}(1-\theta)^{b_0 - 1}}_{\text{\lr{prior}}}
= \theta^{a_0 + y - 1}(1-\theta)^{b_0 + n - y - 1},
\]

پس پسین دوباره یک توزیع \lr{Beta} است:

\[
\theta \mid y \sim \mathrm{Beta}(a_0 + y,\; b_0 + n - y).
\]

وقتی پیشین و پسین از یک خانواده باشند، پیشین را \textbf{\lr{conjugate}}
می‌نامند.

\begin{itemize}
\tightlist
\item
  \(a_0\) و \(b_0\) مثل \textbf{شمارش‌های کاذب} \lr{(pseudo-counts)}
  موفقیت‌ها و شکست‌های پیشین عمل می‌کنند.
\item
  میانگین پسین: \(\dfrac{a_0 + y}{a_0 + b_0 + n}\)؛ میانگین وزنی میانگین
  پیشین \(\frac{a_0}{a_0 + b_0}\) و نسبت مشاهده‌شده‌ی \(y/n\)، با وزن‌های
  \(a_0 + b_0\) و \(n\).
\item
  \textbf{ارتباط با هفته‌های 2 و 4:} تصحیح \lr{log-linear}
  \((y + 0.5)/(n + 1)\) دقیقاً میانگین پسین تحت پیشین
  \(\mathrm{Beta}(0.5, 0.5)\)، به نام \emph{پیشین \lr{Jeffreys}}، است.
  آن «تصحیح» یک برآورد بیزی در لباس مبدل بود.
\end{itemize}

\paragraph{خواندن ریاضی}

\begin{itemize}
\tightlist
\item
  \textbf{توزیع \lr{Beta}} \(\mathrm{Beta}(a, b)\) یک منحنی زنگوله‌مانند
  انعطاف‌پذیر روی بازه‌ی 0 تا 1 است و برای توصیف باور درباره‌ی یک احتمال
  ایده‌آل است. میانگین آن \(a/(a+b)\) است. \(\mathrm{Beta}(1,1)\) صاف
  است؛ \(\mathrm{Beta}(20, 5)\) یک برآمدگی باریک حول 0.8 است.
\item
  قاعده‌ی به‌روزرسانی
  \(\mathrm{Beta}(a_0,\allowbreak b_0) \to \mathrm{Beta}(a_0 + y,\allowbreak b_0 + n - y)\)
  فقط \textbf{شمردن} است: موفقیت‌های مشاهده‌شده را به \(a\) و شکست‌ها را به
  \(b\) اضافه کن. \lr{prior} طوری رفتار می‌کند که انگار قبلاً \(a_0\)
  موفقیت و \(b_0\) شکست دیده‌ایم (شمارش‌های ساختگی یا \lr{pseudo-counts).}
\item
  \textbf{مثال:} \lr{prior} برابر \(\mathrm{Beta}(2, 2)\) («انگار» 2
  یادآوری و 2 شکست)، داده 7 از 10 \ensuremath{\leftarrow} \(\mathrm{Beta}(9, 5)\) با میانگین
  \(9/14 = 0.643\). \lr{prior} مقدار مشاهده‌شده‌ی 0.70 را کمی به سمت 0.5
  می‌کشد، چون 4 مشاهده‌ی ساختگی به 10 مشاهده‌ی واقعی اضافه می‌کند.
\end{itemize}

\subsubsection{\texorpdfstring{4. تقریب شبکه‌ای \lr{(grid
approximation)}}{4. تقریب شبکه‌ای (grid approximation)}}

بیشتر مدل‌های شناختی پیشین \lr{conjugate} ندارند. اگر فقط یک یا دو
پارامتر باشد، می‌توانیم پسین را روی یک شبکه تقریب بزنیم:

\begin{enumerate}
\def\labelenumi{\arabic{enumi}.}
\tightlist
\item
  شبکه‌ای از مقادیر پارامتر
  \(\boldsymbol\theta_1,\allowbreak \dots,\allowbreak \boldsymbol\theta_G\) تعریف کنید که همه‌ی
  مقادیر معقول را بپوشاند.
\item
  در هر نقطه‌ی شبکه \textbf{لگاریتم} پیشین به‌علاوه‌ی \textbf{لگاریتم}
  درست‌نمایی را حساب کنید (لگاریتم‌ها، مثل هفته‌ی 4، از سرریز جلوگیری
  می‌کنند).
\item
  بیشینه را کم کنید، نمایی بگیرید و نرمال کنید تا جمع مقادیر 1 شود (یا
  انتگرالشان 1 شود).
\end{enumerate}

نتیجه یک تقریب گسسته از پسین است که هر خلاصه‌ای را می‌شود از آن حساب کرد.
تقریب شبکه‌ای تا حد دقت شبکه دقیق است و برای ساختن شهود ایده‌آل است. فراتر
از 2 تا 3 پارامتر شکست می‌خورد (نفرین ابعاد، هفته‌ی 3)، و همین انگیزه‌ی
\lr{MCMC} است (هفته‌های 11\lr{--}13).

\subsubsection{5. خلاصه کردن یک پسین}

\begin{longtable}[]{@{}
  >{\raggedright\arraybackslash}p{(\linewidth - 4\tabcolsep) * \real{0.3333}}
  >{\raggedright\arraybackslash}p{(\linewidth - 4\tabcolsep) * \real{0.3333}}
  >{\raggedright\arraybackslash}p{(\linewidth - 4\tabcolsep) * \real{0.3333}}@{}}
\caption{\centering 5. خلاصه کردن یک پسین}\tabularnewline
\toprule\noalign{}
\begin{minipage}[b]{\linewidth}\raggedright
خلاصه
\end{minipage} & \begin{minipage}[b]{\linewidth}\raggedright
تعریف
\end{minipage} & \begin{minipage}[b]{\linewidth}\raggedright
توضیح
\end{minipage} \\
\midrule\noalign{}
\endfirsthead
\toprule\noalign{}
\begin{minipage}[b]{\linewidth}\raggedright
خلاصه
\end{minipage} & \begin{minipage}[b]{\linewidth}\raggedright
تعریف
\end{minipage} & \begin{minipage}[b]{\linewidth}\raggedright
توضیح
\end{minipage} \\
\midrule\noalign{}
\endhead
\bottomrule\noalign{}
\endlastfoot
میانگین پسین & \(\mathbb{E}[\theta \mid \mathbf{y}]\) & خطای مربع مورد
انتظار را کمینه می‌کند \\
میانه‌ی پسین & صدک 50 & در برابر چولگی مقاوم است \\
\lr{MAP} & \(\arg\max_\theta p(\theta\mid\mathbf{y})\) & تحت پیشین تخت
با \lr{MLE} برابر است \\
انحراف معیار پسین & \(\sqrt{\mathrm{Var}[\theta\mid\mathbf{y}]}\) &
همتای بیزی \lr{SE} \\
بازه‌ی دم‌برابر 95\% & صدک 2.5 تا 97.5 & محاسبه‌اش آسان است \\
\textbf{\lr{HDI}} 95\% (بازه‌ی بیشترین چگالی) & باریک‌ترین بازه‌ای که 95\%
پسین را در بر دارد؛ هر مقدار داخل آن از هر مقدار بیرونش محتمل‌تر است &
برای پسین‌های چوله ترجیح دارد \\
\end{longtable}

\textbf{\lr{credible interval} در برابر بازه‌ی اطمینان.} \lr{credible
interval} 95\% \emph{واقعاً} یعنی «با توجه به مدل، پیشین و داده، پارامتر
با احتمال 0.95 در این بازه است». بازه‌ی اطمینان گزاره‌ای درباره‌ی رفتار
بلندمدت یک رویه است (هفته‌ی 5). با پیشین‌های ضعیف و داده‌ی کافی، این دو
اغلب از نظر عددی یکی می‌شوند، ولی تفسیرشان فرق دارد.

\textbf{احتمال پسین فرضیه‌ها} مجانی به دست می‌آید:
\(P(d' > 1 \mid \mathbf{y})\) فقط جرم پسین بالای 1 است.

\subsubsection{\texorpdfstring{6. بررسی‌های پیش‌بینی پیشین \lr{(prior
predictive
checks)}}{6. بررسی‌های پیش‌بینی پیشین (prior predictive checks)}}

\textbf{توزیع پیش‌بینی پیشین} داده‌هایی را توصیف می‌کند که مدل \emph{پیش
از} دیدن هر داده‌ای انتظار دارد:

\[
p(\mathbf{y}^{\text{\lr{rep}}}) = \int p(\mathbf{y}^{\text{\lr{rep}}}\mid\boldsymbol\theta)\, p(\boldsymbol\theta)\, d\boldsymbol\theta.
\]

برای شبیه‌سازی از آن: \(\boldsymbol\theta\) را از پیشین بکشید و بعد از
مدل داده شبیه‌سازی کنید. اگر داده‌های شبیه‌سازی‌شده پوچ به نظر برسند (مثلاً
تقریباً همه‌ی نرخ‌های \lr{hit} دقیقاً 0 یا 1، یا یادآوری‌ای که طی یک ماه
\emph{بیشتر} می‌شود)، پیشین نامعقول است. بررسی‌های پیش‌بینی پیشین نسخه‌ی
بیزی عادت «پیش از برازش شبیه‌سازی کن» هفته‌ی 2 هستند.

\begin{quote}
\textbf{پیشرفته: پسین به‌عنوان پیشین به‌روزشده.} چون برای داده‌های مستقل
\(p(\boldsymbol\theta\mid\mathbf{y}_1,\allowbreak \mathbf{y}_2) \propto p(\mathbf{y}_2\mid\boldsymbol\theta)\, p(\boldsymbol\theta\mid\mathbf{y}_1)\)،
پسین امروز پیشین فرداست. پس به‌روزرسانی بیزی ذاتاً ترتیبی است، و به همین
دلیل استنباط بیزی یک مدل هنجاری محبوب برای \emph{یادگیری} است (هفته‌ی
15).
\end{quote}

\paragraph{برای یک مهندس}

\begin{itemize}
\tightlist
\item
  این همان \textbf{تخمین بیزی} در نظریه‌ی آشکارسازی و تخمین است:
  \lr{posterior} چیزی است که تخمین‌گر \lr{MAP} یا \lr{MMSE} با آن کار
  می‌کند. میانگین \lr{posterior} همان \textbf{تخمین \lr{MMSE}} است؛ مد
  \lr{posterior} \textbf{تخمین \lr{MAP}} است؛ با \lr{prior} صاف،
  \lr{MAP} = \lr{ML.}
\item
  \textbf{به‌روزرسانی \lr{conjugate} بتا\lr{--}دوجمله‌ای} همان ایده‌ی به‌روزرسانی
  \lr{conjugate} گاوسی است، یعنی هسته‌ی اسکالر \textbf{فیلتر کالمن}
  (میانگین و واریانس \lr{prior} + اندازه‌گیری \ensuremath{\leftarrow} میانگین و واریانس
  \lr{posterior)}، فقط اینجا برای یک احتمال به جای یک حالت گاوسی.
\item
  \textbf{به‌روزرسانی ترتیبی} \lr{(«posterior} امروز \lr{prior} فرداست»)
  دقیقاً همان فیلترینگ بیزی بازگشتی است.
\item
  تفاوت در روان‌شناسی: \lr{prior} معمولاً \textbf{ضعیف} است و باید
  \emph{برای خواننده توجیه شود}، و خروجی اصلی یک تصمیم نیست، بلکه توصیف
  عدم قطعیت درباره‌ی یک پارامتر روان‌شناختی است. بررسی پیش‌بینی پیشین
  \lr{(prior predictive check)} عادتی است که \lr{prior}ها را قابل دفاع
  می‌کند.
\end{itemize}

\paragraph{برای مدرس}

\begin{itemize}
\tightlist
\item
  \textbf{\lr{grid} چهارنقطه‌ای را روی تخته انجام دهید} (مثال حل‌شده، بخش
  1) با ستون‌هایی برای \lr{prior}، درست‌نمایی، حاصل‌ضرب و \lr{posterior.}
  کل استنباط بیزی در همین یک جدول است. بعد \lr{prior} را به (0.1، 0.1،
  0.4، 0.4) تغییر دهید و زنده دوباره حساب کنید: دانشجوها اثر \lr{prior}
  را می‌بینند.
\item
  \textbf{از داستان شمارش‌های ساختگی \lr{Beta} استفاده کنید.} \lr{«prior}
  برابر \lr{Beta(2,} 2) مثل این است که قبلاً 2 یادآوری و 2 شکست دیده
  باشیم.» دانشجوهایی که \lr{prior} را دلبخواهی می‌دانند، وقتی آن را داده‌ی
  خیالی ببینند آرام می‌شوند.
\item
  \textbf{مستقیماً به نگرانی «ذهنی بودن» جواب دهید.} هر تحلیلی فرض دارد
  (درست‌نمایی هم یک فرض است)؛ تحلیل بیزی یکی دیگر از فرض‌ها را صریح می‌کند،
  و تحلیل حساسیت (هفته‌ی 10) بررسی می‌کند که چقدر مهم است.
\item
  \textbf{سؤال‌های رایج.}

  \begin{itemize}
  \tightlist
  \item
    از \lr{prior}های \lr{weakly informative} استفاده کنید، آن‌ها را با
    شبیه‌سازی پیش‌بینی پیشین بررسی کنید و حساسیت را گزارش کنید.
  \item
    \emph{«آیا \lr{credible interval} همان بازه‌ی اطمینان است؟»} با
    \lr{prior} ضعیف و داده‌ی کافی اغلب از نظر عددی شبیه‌اند، ولی فقط
    \lr{credible interval} یعنی «با احتمال 95\% پارامتر اینجاست».
  \item
    \emph{«پس مغز بیزی است؟»} این سؤال دیگری است (هفته‌ی 15). این هفته
    بیز ابزاری برای پژوهشگر است.
  \end{itemize}
\end{itemize}

\begin{center}\rule{0.5\linewidth}{0.5pt}\end{center}

\Needspace{11\baselineskip}

\subsection{مثال حل‌شده}

\subsubsection{بخش 1: یک شبکه‌ی چهارنقطه‌ای با دست}

یک شرکت‌کننده \(y = 7\) از \(n = 10\) آیتم را به یاد می‌آورد.

\begin{longtable}[]{@{}
  >{\raggedright\arraybackslash}p{(\linewidth - 8\tabcolsep) * \real{0.2000}}
  >{\raggedright\arraybackslash}p{(\linewidth - 8\tabcolsep) * \real{0.2000}}
  >{\raggedright\arraybackslash}p{(\linewidth - 8\tabcolsep) * \real{0.2000}}
  >{\raggedright\arraybackslash}p{(\linewidth - 8\tabcolsep) * \real{0.2000}}
  >{\raggedright\arraybackslash}p{(\linewidth - 8\tabcolsep) * \real{0.2000}}@{}}
\caption{\centering بخش 1: یک شبکه‌ی چهارنقطه‌ای با دست}\tabularnewline
\toprule\noalign{}
\begin{minipage}[b]{\linewidth}\raggedright
\(\theta\)
\end{minipage} & \begin{minipage}[b]{\linewidth}\raggedright
پیشین
\end{minipage} & \begin{minipage}[b]{\linewidth}\raggedright
درست‌نمایی \(\theta^7(1-\theta)^3\) (بدون ثابت)
\end{minipage} & \begin{minipage}[b]{\linewidth}\raggedright
پیشین × درست‌نمایی
\end{minipage} & \begin{minipage}[b]{\linewidth}\raggedright
پسین
\end{minipage} \\
\midrule\noalign{}
\endfirsthead
\toprule\noalign{}
\begin{minipage}[b]{\linewidth}\raggedright
\(\theta\)
\end{minipage} & \begin{minipage}[b]{\linewidth}\raggedright
پیشین
\end{minipage} & \begin{minipage}[b]{\linewidth}\raggedright
درست‌نمایی \(\theta^7(1-\theta)^3\) (بدون ثابت)
\end{minipage} & \begin{minipage}[b]{\linewidth}\raggedright
پیشین × درست‌نمایی
\end{minipage} & \begin{minipage}[b]{\linewidth}\raggedright
پسین
\end{minipage} \\
\midrule\noalign{}
\endhead
\bottomrule\noalign{}
\endlastfoot
0.2 & 0.25 & 0.0000066 & 0.0000016 & \textbf{0.002} \\
0.4 & 0.25 & 0.0003539 & 0.0000885 & \textbf{0.092} \\
0.6 & 0.25 & 0.0017916 & 0.0004479 & \textbf{0.468} \\
0.8 & 0.25 & 0.0016777 & 0.0004194 & \textbf{0.438} \\
جمع & 1 & & 0.0009574 & 1.000 \\
\end{longtable}

هر مقدار پسین = (پیشین × درست‌نمایی) / جمع. ضریب دوجمله‌ای حذف شد چون در
نرمال‌سازی ساده می‌شود. تفسیر: بعد از داده، 0.6 و 0.8 تقریباً به یک اندازه
محتمل‌اند و 0.2 عملاً کنار گذاشته شده است.

\subsubsection{\texorpdfstring{بخش 2: همان داده با یک پیشین \lr{Beta}
پیوسته}{بخش 2: همان داده با یک پیشین Beta پیوسته}}

پیشین \(\mathrm{Beta}(2, 2)\) (میانگین 0.5، با ترجیح ملایم برای مقادیر
میانی). پسین: \(\mathrm{Beta}(2 + 7,\allowbreak 2 + 3) = \mathrm{Beta}(9,\allowbreak 5)\).

\begin{itemize}
\tightlist
\item
  میانگین پسین \(= 9/14 = 0.643\): بین میانگین پیشین (0.5) و نسبت
  مشاهده‌شده (0.7)، نزدیک‌تر به داده چون \(n = 10 > a_0 + b_0 = 4\).
\item
  \lr{MAP} (مُد) \(= (9 - 1)/(9 + 5 - 2) = 0.667\).
\item
  \lr{credible interval} دم‌برابر 95\% \(= [0.386, 0.861]\).
\end{itemize}

\textbf{تفسیر.} با یک پیشین ملایم و 10 آیتم، احتمال یادآوری احتمالاً حدود
0.64 است، ولی هر مقداری از حدود 0.39 تا 0.86 هنوز \lr{credible} است. ده
آیتم چیز زیادی به ما نمی‌گوید.

\subsubsection{بخش 3: با داده‌ی بیشتر پیشین کمتر مهم است}

با \(y = 30\) از \(n = 50\):

\begin{longtable}[]{@{}
  >{\raggedright\arraybackslash}p{(\linewidth - 6\tabcolsep) * \real{0.2500}}
  >{\raggedright\arraybackslash}p{(\linewidth - 6\tabcolsep) * \real{0.2500}}
  >{\raggedright\arraybackslash}p{(\linewidth - 6\tabcolsep) * \real{0.2500}}
  >{\raggedright\arraybackslash}p{(\linewidth - 6\tabcolsep) * \real{0.2500}}@{}}
\caption{\centering بخش 3: با داده‌ی بیشتر پیشین کمتر مهم است}\tabularnewline
\toprule\noalign{}
\begin{minipage}[b]{\linewidth}\raggedright
پیشین
\end{minipage} & \begin{minipage}[b]{\linewidth}\raggedright
پسین
\end{minipage} & \begin{minipage}[b]{\linewidth}\raggedright
میانگین
\end{minipage} & \begin{minipage}[b]{\linewidth}\raggedright
\lr{credible interval} 95\%
\end{minipage} \\
\midrule\noalign{}
\endfirsthead
\toprule\noalign{}
\begin{minipage}[b]{\linewidth}\raggedright
پیشین
\end{minipage} & \begin{minipage}[b]{\linewidth}\raggedright
پسین
\end{minipage} & \begin{minipage}[b]{\linewidth}\raggedright
میانگین
\end{minipage} & \begin{minipage}[b]{\linewidth}\raggedright
\lr{credible interval} 95\%
\end{minipage} \\
\midrule\noalign{}
\endhead
\bottomrule\noalign{}
\endlastfoot
\(\mathrm{Beta}(1,1)\) (تخت) & \(\mathrm{Beta}(31, 21)\) & 0.596 &
{[}0.461, 0.724{]} \\
\(\mathrm{Beta}(0.5, 0.5)\) \lr{(Jeffreys)} &
\(\mathrm{Beta}(30.5, 20.5)\) & 0.598 & {[}0.462, 0.727{]} \\
\(\mathrm{Beta}(20, 5)\) (آگاهی‌بخش، «\ensuremath{\approx} 0.8») & \(\mathrm{Beta}(50, 25)\)
& 0.667 & {[}0.557, 0.768{]} \\
\end{longtable}

پیشین‌های تخت و \lr{Jeffreys} عملاً یک جواب می‌دهند که با بازه‌ی \lr{profile
likelihood} هفته‌ی 5، {[}0.462, 0.728{]}، هم برابر است. پیشین آگاهی‌بخش،
که ارزشش 25 مشاهده‌ی کاذب است، هنوز برآورد را به‌طور قابل توجهی به سمت 0.8
می‌کشد: 50 مشاهده برای غلبه بر یک پیشین قوی کافی نیست. اینکه این خوب است
یا نه به درست بودن دانش پیشین بستگی دارد.

\begin{center}\rule{0.5\linewidth}{0.5pt}\end{center}

\Needspace{11\baselineskip}

\subsection{\texorpdfstring{\lr{lab} (2 ساعت)}{lab (2 ساعت)}}

\subsubsection{اهداف}

\begin{enumerate}
\def\labelenumi{\arabic{enumi}.}
\tightlist
\item
  محاسبه‌ی پسین‌های \lr{beta-binomial} به‌صورت تحلیلی و با تقریب شبکه‌ای، و
  دیدن مصالحه‌ی پیشین و داده.
\item
  تماشای تیزتر شدن پسین با داده (به‌روزرسانی ترتیبی).
\item
  محاسبه‌ی پسین شبکه‌ای دوبعدی برای \lr{SDT} (\(d'\)، \(c\))؛ خلاصه کردن
  آن با میانگین‌ها، \lr{HDI}ها و احتمال‌های پسین.
\item
  استفاده از شبیه‌سازی پیش‌بینی پیشین برای افشای یک پیشین «مبهم» بد.
\end{enumerate}

\textbf{زمان‌بندی:} بخش \lr{A} 25 دقیقه · بخش \lr{B} 15 دقیقه · بخش
\lr{C} 40 دقیقه · بخش \lr{D} 25 دقیقه · بخش \lr{E} 15 دقیقه.

\subsubsection{\texorpdfstring{بخش \lr{A: beta-binomial}، تحلیلی و روی
شبکه (25 دقیقه)}{بخش A: beta-binomial، تحلیلی و روی شبکه (25 دقیقه)}}

پسین‌های شبکه‌ای و دقیق بر هم منطبق‌اند. ببینید هر پسین نسبت به پیشین و
درست‌نمایی خودش کجا قرار می‌گیرد.

\subsubsection{\texorpdfstring{بخش \lr{B:} به‌روزرسانی ترتیبی (15
دقیقه)}{بخش B: به‌روزرسانی ترتیبی (15 دقیقه)}}

\textbf{به این دقت کنید:} پسین از خود پیشین شروع می‌شود، با جمع شدن
آیتم‌ها باریک می‌شود و حول مقدار واقعی متمرکز می‌شود. پهنایش تقریباً مثل
\(1/\sqrt{n}\) کم می‌شود، همان‌طور که \lr{SE} در هفته‌ی 5 کم می‌شد.

\subsubsection{\texorpdfstring{بخش \lr{C:} یک پسین شبکه‌ای دوپارامتری
برای \lr{SDT} (40
دقیقه)}{بخش C: یک پسین شبکه‌ای دوپارامتری برای SDT (40 دقیقه)}}

شرکت‌کننده‌ی هفته‌ی 2: 68 \lr{hit} از 100 \lr{trial} علامت و 11 \lr{false
alarm} از 100 \lr{trial} نویز. پیشین‌ها: \(d' \sim \mathcal{N}(1, 1.5)\)
(حساسیت معمولاً مثبت و به‌ندرت بالای 4 است) و \(c \sim \mathcal{N}(0, 1)\)
(سوگیری معمولاً متوسط است).

آن را همین‌طور گزارش کنید، نه به‌عنوان قطعیت.)

میانگین پسین و \lr{MAP} در مقایسه با \lr{MLE} چطورند؟ چرا اینجا این‌قدر
نزدیک‌اند؟ 2. \(P(c > 0 \mid \text{\lr{data}})\) 3. انحراف معیار پیشین \(d'\)
را به 10 افزایش دهید، بعد به 0.3 کم کنید (با مرکز 1).

\subsubsection{\texorpdfstring{بخش \lr{D:} بررسی‌های پیش‌بینی پیشین (25
دقیقه)}{بخش D: بررسی‌های پیش‌بینی پیشین (25 دقیقه)}}

پیش از اعتماد به یک پیشین، به داده‌هایی که نتیجه می‌دهد نگاه کنید. پیشین
ما را با یک پیشین «مبهم»، \(d' \sim \mathcal{N}(0, 10)\) و
\(c \sim \mathcal{N}(0, 10)\)، مقایسه کنید که گاهی به این دلیل استفاده
می‌شود که «بی‌اطلاع» است.

\textbf{به این دقت کنید:} پیشین «مبهم» اصلاً خنثی نیست. روی مقیاس احتمال،
96\% مجموعه‌داده‌های شبیه‌سازی‌شده را در نرخ \lr{hit} یا \lr{false alarm}
دقیقاً 0 یا 1 قرار می‌دهد؛ چیزی که در آزمایش‌های تشخیص واقعی عملاً هرگز
اتفاق نمی‌افتد. این پیشین باوری قوی و پوچ درباره‌ی داده بیان می‌کند. ابهام
روی مقیاس پارامتر، ابهام روی مقیاس داده نیست.

پیشین خود ما هم کامل نیست: حدود یک‌چهارم مجموعه‌داده‌های شبیه‌سازی‌شده‌اش
\(H < F\) دارند، چون \(\mathcal{N}(1, 1.5)\) حدود 25\% جرمش را روی
\(d'\) منفی می‌گذارد. برای تکلیفی که می‌دانیم شرکت‌کننده‌ها بالاتر از شانس
عمل می‌کنند، پیشینی محدود به \(d' > 0\) (مثلاً نیم‌نرمال یا \lr{Gamma)}
قابل دفاع‌تر است. امتحانش کنید و ببینید آیا پسین این شرکت‌کننده را تغییر
می‌دهد. (راهنمایی: با 200 \lr{trial} آگاهی‌بخش، تقریباً نه.)

\subsubsection{\texorpdfstring{بخش \lr{E:} تفسیر (15
دقیقه)}{بخش E: تفسیر (15 دقیقه)}}

یک «پاراگراف نتایج» کوتاه (4 تا 6 جمله) برای این شرکت‌کننده بنویسید، طوری
که در یک مقاله می‌آید. پیشین‌ها، میانگین‌های پسین و \lr{HDI}های 95\% برای
\(d'\) و \(c\)، \(P(d' > 1 \mid \text{\lr{data}})\)، و یک جمله درباره‌ی حساسیت
به پیشین (از سؤال 3 بخش \lr{C)} را گزارش کنید.

\subsubsection{تمرین‌ها (به ترتیب سختی)}

\textbf{تمرین 1 \ensuremath{\star}.} به‌صورت جبری نشان دهید که میانگین پسین مدل
\lr{beta-binomial} میانگین وزنی میانگین پیشین و \(y/n\) است. وزن‌ها
چیستند؟

\textbf{تمرین 2 \ensuremath{\star}.} با شبکه‌ی بخش \lr{A}، \lr{HDI} 95\% و بازه‌ی دم‌برابر
را برای \(\mathrm{Beta}(1 + 2, 1 + 18)\) (2 موفقیت از 20) حساب کنید. چرا
با هم فرق دارند؟

\textbf{تمرین 3 \ensuremath{\star}\ensuremath{\star} (فراموشی با شبکه).} پسین شبکه‌ای \((A, d)\) را برای
مدل فراموشی توانی با داده‌های هفته‌ی 2، با پیشین‌های
\(A \sim \mathrm{Beta}(8, 2)\) و
\(d \sim \mathrm{Gamma}(\text{\lr{shape }} 2, \text{\lr{scale }} 0.3)\) حساب کنید.
میانگین‌های پسین و \lr{HDI}های 95\% را با \lr{MLE} و بازه‌های \lr{profile}
و \lr{bootstrap} هفته‌ی 5 مقایسه کنید.

\textbf{تمرین 4 \ensuremath{\star}\ensuremath{\star} (هم‌ارزی \lr{MAP} و \lr{MLE).}} در بخش \lr{C} هر دو
پیشین را با پیشین‌های تخت روی شبکه جایگزین کنید. نشان دهید \lr{MAP} برابر
\lr{MLE} است. چرا \emph{میانگین} پسین هنوز کمی فرق دارد؟

\textbf{تمرین 5 \ensuremath{\star}\ensuremath{\star}\ensuremath{\star} (پیشرفته: ترتیبی در برابر یکجا).} برای مدل \lr{SDT}،
پسین شبکه‌ای را \lr{trial} به \lr{trial} (پسین هر \lr{trial} پیشین
\lr{trial} بعدی است) برای 200 \lr{trial} شبیه‌سازی‌شده به‌روز کنید و نشان
دهید پسین نهایی (تا حد خطای عددی) با پسین یکجای همه‌ی \lr{trial}ها یکی
است. انحراف معیار پسین \(d'\) را در برابر تعداد \lr{trial}ها رسم کنید.

\subsubsection{\texorpdfstring{\lr{lab check-off}}{lab check-off}}

نشان دهید: (1) نمودار سه‌قسمتی پیشین/درست‌نمایی/پسین از بخش \lr{A}، (2)
جدول خلاصه‌ی \lr{SDT} و پسین مشترک از بخش \lr{C}، و (3) دو نمودار پیش‌بینی
پیشین از بخش \lr{D} با یک جمله درباره‌ی اینکه پیشین مبهم چه مشکلی دارد.

\begin{center}\rule{0.5\linewidth}{0.5pt}\end{center}

\subsection{\texorpdfstring{ارتباط با هوش مصنوعی و
\lr{LLM}ها}{ارتباط با هوش مصنوعی و LLMها}}

\textbf{مدل واقعاً چقدر خوب است؟ یک جواب بیزی برای ارزیابی‌های کوچک.}
پژوهشگرها اغلب یک \lr{LLM} را روی مجموعه‌ی کوچکی از آیتم‌ها، مثلاً 25
مسئله‌ی دست‌نویس، می‌آزمایند. مدل \lr{beta-binomial} این هفته همان چیزی را
می‌دهد که چنین مطالعه‌هایی معمولاً ندارند: یک پسین کامل برای دقت، گزاره‌های
احتمالی مستقیم («احتمال اینکه دقت از 60\% بیشتر باشد») و مقایسه‌ی مستقیم
دو مدل.

با 25 آیتم، پهنای بازه‌های پسین حدود 0.35 است، و احتمال پسین 19\% وجود
دارد که در واقع \lr{B} مدل بهتر باشد.

نتیجه‌گیری‌ها چقدر عوض می‌شوند؟ - \ensuremath{\star}\ensuremath{\star} مقایسه‌ی بالا دو مدل را مستقل در نظر
می‌گیرد، با اینکه به همان آیتم‌ها جواب داده‌اند. چرا این مشکل است (کادر
هفته‌ی 5 را ببینید)؟ طرح کلی مدلی را بکشید که جفت بودن را رعایت کند. - \ensuremath{\star}\ensuremath{\star}
بحث: برای دقت یک مدل کاملاً تازه روی یک آزمون چهارگزینه‌ای چه پیشینی
استفاده می‌کنید، و چرا نباید جرم زیادی زیر 0.25 بگذارد؟

\begin{center}\rule{0.5\linewidth}{0.5pt}\end{center}

\subsection{اشتباهات رایج و برداشت‌های نادرست}

\begin{enumerate}
\def\labelenumi{\arabic{enumi}.}
\tightlist
\item
  \textbf{«پیشین استنباط بیزی را ذهنی و نتایج را دلبخواهی می‌کند.»}
  پیشین‌ها فرض‌های صریح‌اند، درست مثل انتخاب درست‌نمایی. گزارششان کنید،
  توجیهشان کنید و حساسیت را بررسی کنید. با داده‌ی کافی، پیشین‌های معقول به
  نتیجه‌گیری‌های مشابه می‌رسند.
\item
  \textbf{«پیشین‌های تخت یا خیلی پهن بی‌اطلاع‌اند.»} می‌توانند روی مقیاس
  داده به‌شدت آگاهی‌بخش باشند (بخش \lr{D).} از پیشین‌های ضعیف‌آگاهی‌بخش
  استفاده کنید و با شبیه‌سازی پیش‌بینی پیشین بررسی‌شان کنید.
\item
  \textbf{انتخاب پیشین بعد از نگاه کردن به داده.} این یعنی دو بار
  استفاده از داده.
\item
  \textbf{ضرب احتمال‌ها روی شبکه بدون لگاریتم،} که سرریز می‌کند (هفته‌ی 4).
  با لگاریتم پیشین + لگاریتم درست‌نمایی کار کنید و پیش از نمایی گرفتن
  بیشینه را کم کنید.
\item
  \textbf{شبکه‌ی خیلی درشت یا خیلی باریک.} اگر پسین در لبه‌ی شبکه بریده
  شود، همه‌ی خلاصه‌ها غلط‌اند. بررسی کنید که پسین در لبه‌ها نزدیک صفر باشد.
\item
  \textbf{خواندن \lr{credible interval} به جای بازه‌ی اطمینان یا برعکس.}
  فقط \lr{credible interval} جمله‌ی «با احتمال 95\% پارامتر اینجاست» (با
  توجه به مدل و پیشین) را پشتیبانی می‌کند.
\item
  \textbf{فقط گزارش \lr{MAP.}} نکته‌ی پسین کل توزیع است. یک مقدار مرکزی
  \emph{و} یک بازه، و احتمال فرضیه‌هایی را که برایتان مهم‌اند گزارش کنید.
\end{enumerate}

\begin{center}\rule{0.5\linewidth}{0.5pt}\end{center}

\subsection{خودآزمایی}

\textbf{1.} قاعده‌ی بیز را برای پارامتر \(\theta\) بنویسید و هر جمله را
نام ببرید. \textbf{پاسخ}
\(p(\theta\mid y) = p(y\mid\theta)p(\theta)/p(y)\): پسین = درست‌نمایی ×
پیشین / شواهد \lr{(marginal likelihood).} برای برآورد،
\(p(\theta\mid y) \propto p(y\mid\theta)p(\theta)\).

\textbf{2.} پیشین \(\mathrm{Beta}(3, 3)\)، داده: 12 یادآوری از 20. پسین
و میانگینش چیست؟ \textbf{پاسخ}
\(\mathrm{Beta}(3 + 12,\allowbreak 3 + 8) = \mathrm{Beta}(15,\allowbreak 11)\)، با میانگین
\(15/26 = 0.577\).

\textbf{3.} چرا با افزایش \(n\) پسین به \lr{MLE} نزدیک‌تر می‌شود؟
\textbf{پاسخ} لگاریتم درست‌نمایی مجموعی روی مشاهدات است و با \(n\) بزرگ
می‌شود، در حالی که لگاریتم پیشین ثابت می‌ماند. با داده‌ی زیاد درست‌نمایی بر
حاصل‌ضرب غالب می‌شود و پسین حول بیشینه‌ی درست‌نمایی متمرکز می‌شود.

\textbf{4.} تفاوت \lr{HDI} 95\% و بازه‌ی دم‌برابر 95\% چیست؟ کی بر هم
منطبق‌اند؟ \textbf{پاسخ} \lr{HDI} باریک‌ترین بازه‌ای است که 95\% جرم پسین
را در بر دارد و هر مقدار داخلش چگالی بیشتری از هر مقدار بیرونش دارد.
بازه‌ی دم‌برابر از هر دم 2.5\% را می‌بُرد. برای پسین‌های متقارن تک‌مُدی بر هم
منطبق‌اند و برای پسین‌های چوله فرق دارند.

\textbf{5.} چرا تقریب شبکه‌ای برای مدل‌هایی با پارامترهای زیاد غیرعملی
می‌شود؟ \textbf{پاسخ} تعداد نقاط شبکه با تعداد پارامترها به‌طور نمایی رشد
می‌کند (مثلاً 100 نقطه برای هر پارامتر یعنی \(100^k\) ارزیابی). این نفرین
ابعاد است. نمونه‌گیری \lr{MCMC} (هفته‌های 11\lr{--}13) از آن دوری می‌کند.

\textbf{6.} همکاری «برای احتیاط» از \(d' \sim \mathcal{N}(0, 100)\)
استفاده می‌کند. برای قانع کردنش که این ایده‌ی بدی است چه می‌کنید؟
\textbf{پاسخ} یک شبیه‌سازی پیش‌بینی پیشین: نشان دهید این پیشین نرخ‌های
\lr{hit} و \lr{false alarm}ی نتیجه می‌دهد که تقریباً همیشه دقیقاً 0 یا 1
هستند؛ باوری پوچ درباره‌ی داده‌ی تشخیص واقعی. به‌جایش یک پیشین ضعیف‌آگاهی‌بخش
پیشنهاد کنید.

\textbf{7.} این جمله را توضیح دهید: «تصحیح \lr{log-linear} برای نرخ‌های
حدی یک برآورد بیزی است». \textbf{پاسخ} نرخ تصحیح‌شده‌ی
\((y + 0.5)/(n + 1)\) برابر میانگین پسین \(\theta\) تحت پیشین
\(\mathrm{Beta}(0.5, 0.5)\) است: \((0.5 + y)/(0.5 + 0.5 + n)\). این
تصحیح به‌طور ضمنی از هر نوع نیم مشاهده‌ی کاذب اضافه می‌کند.

\begin{center}\rule{0.5\linewidth}{0.5pt}\end{center}

\subsection{تکلیف}

\begin{itemize}
\tightlist
\item
  \textbf{\lr{A3} (مقایسه و بازیابی مدل)} و \textbf{پیشنهاد پروژه} در
  ابتدای درس این هفته تحویل داده می‌شوند.
\item
  این هفته تکلیف نمره‌دار جدیدی نداریم. \textbf{\lr{A4} (تحلیل بیزی با
  شبکه)} در هفته‌ی 10 منتشر می‌شود.
\item
  \textbf{تمرین اختیاری (بدون نمره):} تمرین 3 \lr{lab} (فراموشی با شبکه)
  هسته‌ی \lr{A4} است. اگر الان شروعش کنید در وقتتان صرفه‌جویی می‌شود.
\end{itemize}

\begin{center}\rule{0.5\linewidth}{0.5pt}\end{center}

\subsection{مطالعه‌ی بیشتر}

\begin{itemize}
\tightlist
\item
  \textbf{\lr{Kruschke} (2014)، فصل‌های 5\lr{--}6.} قاعده‌ی بیز، و استنباط یک
  احتمال دوجمله‌ای با تحلیل ریاضی دقیق (مدل \lr{beta-binomial).} خیلی
  ساده و تصویری. \emph{(منبع اصلی.)}
\item
  \textbf{\lr{Lee} \lr{\&} \lr{Wagenmakers} (2014)، فصل‌های 1 و 3.} مبانی
  تحلیل بیزی، و استنباط با توزیع دوجمله‌ای.
\item
  \textbf{\lr{Farrell} \lr{\&} \lr{Lewandowsky} (2018)، فصل 6.} برآورد بیزی
  پارامتر: مفاهیم پایه، مرتبط با چارچوب درست‌نمایی هفته‌ی 4.
\item
  \textbf{عظیمی‌فر و احمدی (1396)}، جلد 1: بخش برآورد بیزی پارامتر و
  ارتباطش با برآورد \lr{ML} {[}شماره‌ی فصل بررسی شود{]}.
\end{itemize}

\emph{پیش‌نمایش هفته‌ی 10:} پسین درباره‌ی پارامترها \emph{درون} یک مدل حرف
می‌زند. هفته‌ی بعد می‌پرسیم داده‌ها باور ما را \emph{بین} مدل‌ها چطور تغییر
می‌دهند، با استفاده از شواهد \(p(\mathbf{y}\mid\mathcal{M})\) که تا حالا
به‌عنوان ثابت نرمال‌ساز نادیده‌اش گرفته‌ایم، و چطور بررسی کنیم که پیش‌بینی‌های
پسین یک مدل شبیه داده‌ی واقعی هستند یا نه.
